\chapter{多主体信息认知传播评估系统的实现与测试}

\section{系统实现}

本节在第三章系统概要设计的基础上，按照六个功能模块逐一阐述具体实现方案，各模块的编号与第3.3.2节一致。

\subsection{数据导入与预处理模块的实现}

该模块实现从原始JSON数据到模型可用特征的完整预处理流水线。预处理完成后将数据写入数据库进行持久化存储，数据库表结构设计详见第3.3.4节。预处理流水线按以下步骤依次执行。

\textbf{（1）格式校验与去重}

系统首先校验导入数据的字段完整性，检查每条记录是否包含全部必要字段；随后基于tweet\_id进行去重，保留首次出现的记录。

\textbf{（2）质量过滤}

质量过滤首先剔除crawl\_time缺失的样本，此类样本无法计算随访时间，对两个评估模块均不可用。对tweet\_text为空的样本，系统保留节点但标记为传播价值评估不可用：传播延续预测模块（IHPE）需要维护评论树完整性，保留空文本节点以近零向量表示；传播价值评估模块（Group-aware Fusion TabTransformer）需要文本特征，空文本节点不参与其推理。系统统一按传播延续预测模块的宽松策略执行质量过滤（仅剔除crawl\_time缺失样本），所有保留样本写入Comment表，并通过标记字段区分各记录对传播价值评估的可用性；传播价值评估模块在读取时仅筛选标记为可用的记录。

\textbf{（3）右删失过滤}

系统计算每条评论的随访时间$A = \text{crawl\_time} - \text{created\_at}$，剔除$A < T$（$T=12\text{h}$，即43200秒）的样本。两个评估模型均在满足$A \geq T$约束的样本上训练；若将$A < T$的评论送入模型推理，其评论树仍处于演化状态（子回复可能尚未产生），入模特征与训练时的数据分布不一致，预测结果不可靠。

\textbf{（4）评论树结构重构}

系统基于in\_reply\_to\_tweet\_id字段建立评论的父子关系，通过广度优先搜索（BFS）遍历计算每个节点的树深度和所属评论树根节点。对未被BFS覆盖的孤立节点（父节点不在当前数据集内），作为独立根处理。

\textbf{（5）特征提取}

特征提取涵盖四个方面：首先，BGE文本编码调用BAAI/bge-base-en-v1.5模型，将评论文本编码为768维稠密向量；其次，情感分析调用cardiffnlp/twitter-roberta-base-sentiment-latest模型，输出3维情感概率向量；此外，外生上下文特征从发布时间中提取小时（0--23）和星期（0--6），从话题地区提取地区信息，均采用独热编码；最后，主体标识直接读取导入数据中的agent\_type字段（0=Human, 1=Agent），该字段由用户在数据准备阶段完成标注。

数据导入与预处理模块的处理流程如图\ref{fig:import_flow}所示，对应run\_import\_pipeline()函数的七步流水线。系统接收JSON文件后创建导入任务记录（标记为处理中），随后依次执行：格式校验（检查11个必要字段是否完整，不通过则终止并标记为失败）、评论ID去重、质量过滤（剔除采集时间缺失样本，标记空文本节点为传播价值评估不可用）、右删失过滤（剔除随访时间$A<T$的样本）、评论树结构重构（BFS计算节点深度与所属根节点）、特征提取（BGE文本编码、情感分析、外生上下文特征）。最后执行数据库写入，跳过已存在的评论ID以支持跨导入去重，批量写入Comment与CommentFeature表，完成后更新各阶段计数并标记为已完成。

\begin{figure}[H]
  \centering
  \includegraphics[width=0.85\textwidth]{G6-1/图6-1 数据导入与预处理模块流程图.cropped.pdf}
  \bicaption{数据导入与预处理模块流程图}{Flowchart of Data Import and Preprocessing Module}
  \label{fig:import_flow}
\end{figure}

\subsection{数据存储模块的实现}

该模块负责系统全流程数据的结构化持久存储，为数据导入模块提供写入接口，为评估模块、对比分析模块和可视化模块提供读取接口。

\textbf{（1）存储引擎与会话管理}

系统采用SQLite作为存储引擎，通过SQLAlchemy ORM框架实现数据库访问。数据库引擎在应用启动时初始化，连接配置适配FastAPI的异步请求处理模式，并通过会话工厂管理数据库会话的创建与回收。数据库会话以FastAPI依赖注入方式提供给各业务模块，在请求结束时自动关闭，确保连接资源的正确释放。

\textbf{（2）ORM表定义}

基于SQLAlchemy的声明式映射，定义五张数据表的ORM模型类，各类的字段类型、约束条件与第3.3.4节的表结构设计一一对应：ImportTask（导入任务表）、Comment（评论数据表）、CommentFeature（评论特征表）、EvaluationResult（评估结果表）、AnalysisReport（对比分析结果表）。

\textbf{（3）表间关系与级联删除}

通过SQLAlchemy的关系映射机制建立表间关联：ImportTask与Comment为一对多关系，ImportTask与AnalysisReport为一对多关系，Comment与CommentFeature为一对一关系，Comment与EvaluationResult为一对一关系。所有外键均配置级联删除策略，确保删除导入任务时自动级联清除关联的评论数据、特征数据、评估结果与分析报告。

\textbf{（4）数据库初始化}

应用启动时通过FastAPI的生命周期管理器触发数据库初始化，根据ORM模型定义自动创建所有尚不存在的数据表，无需手动执行数据库迁移脚本。

数据存储模块的ORM模型类结构如图\ref{fig:orm_class}所示。

\begin{figure}[H]
  \centering
  \includegraphics[width=0.9\textwidth]{G6-2/图6-2 数据存储模块类图.cropped.pdf}
  \bicaption{数据存储模块类图}{Class Diagram of Data Storage Module}
  \label{fig:orm_class}
\end{figure}

\subsection{评论传播价值评估模块的实现}

该模块封装第四章提出的Group-aware Fusion TabTransformer模型，提供评论级传播价值评估服务。

\textbf{（1）模型加载与管理}

系统加载预训练的Group-aware Fusion TabTransformer模型权重，模型配置参数与第四章实验保持一致（隐层维度32，Transformer层数2，注意力头数4）。

\textbf{（2）特征组装}

推理前从预处理结果中提取单条评论的四组特征：G1外生上下文特征包括话题地区独热编码（类别数与训练数据集一致）、发布小时独热编码（24维）和星期独热编码（7维）；G2文本语义特征为768维BGE编码向量；G3情感倾向特征为3维情感概率向量；G4主体标识特征为1维agent\_type。

\textbf{（3）推理与评分}

系统调用第四章Group-aware Fusion TabTransformer模型进行推理，输出传播价值评分（Sigmoid概率值，取值范围0至1），支持批量推理以提高效率。

\textbf{（4）结果输出}

推理完成后为每条评论附加传播价值评分字段，支持按评分降序排序和按阈值筛选（如评分大于0.5的高价值评论）。

\subsection{传播延续预测模块的实现}

该模块封装第五章提出的IHPE模型，提供路径级传播延续预测服务。

\textbf{（1）模型加载与管理}

系统加载预训练的IHPE模型权重，模型配置参数与第五章实验保持一致（隐层维度128，LSTM层数2，dropout率0.2）。

\textbf{（2）路径构造与特征组装}

系统从评论树结构中提取目标节点的完整路径$P_k = [v_0, v_1, \ldots, v_k]$，为路径上每个节点提取特征（BGE文本编码768维、情感概率3维、主体类型嵌入），并为路径上每条边计算交互类型（H$\to$H、H$\to$A、A$\to$H、A$\to$A）。

\textbf{（3）推理与预测}

系统调用第五章IHPE模型进行推理，输出2跳增长概率（取值范围0至1），支持批量路径推理。

\textbf{（4）结果输出}

推理完成后为每个节点附加2跳增长概率字段，支持按概率降序排序以识别高传播潜力的讨论链。

评论传播价值评估模块（第6.1.3节）与传播延续预测模块在系统架构中具有对称的结构设计，图\ref{fig:eval_class}以统一的类图展示两个评估模块的静态结构及其与支撑层、数据存储层的关系。评估推理的公共入口为eval\_service模块的run\_eval()函数，内部分别调用\_run\_tab\_transformer()和\_run\_ihpe()两条推理链路。支撑层的tab\_transformer模块和ihpe模块各自封装一个PyTorch模型实例（GroupFusionTabTransformer、IHPE），通过模块级的infer()/infer\_paths()接口向业务层提供推理能力，模型权重可独立替换而不影响业务逻辑。

特征组装与路径构造逻辑分别内联于\_run\_tab\_transformer()（组装G1外生上下文、G2文本语义、G3情感倾向、G4主体标识四组特征）和\_run\_ihpe()（按深度排序构建根到目标路径并利用父路径缓存计算交互类型编码）中，从数据存储模块读取预提取特征并转换为模型所需的输入格式。

\begin{figure}[H]
  \centering
  \includegraphics[width=0.9\textwidth]{G6-3/图6-3 评估子系统类图.cropped.pdf}
  \bicaption{评估子系统类图}{Class Diagram of Evaluation Subsystem}
  \label{fig:eval_class}
\end{figure}

\subsection{多主体对比分析模块的实现}

该模块聚合评论传播价值评估模块和传播延续预测模块的输出，按主体类型进行分层对比分析，是系统实现"间接刻画信息认知传播特征"这一总任务的核心模块。

评论传播价值评估模块与传播延续预测模块从两个互补视角刻画同一节点的传播潜力：前者（评论级）回答"该评论本身能否吸引直接回复"，反映评论内容的话题吸引力；后者（路径级）回答"该节点所在的讨论链能否持续延伸"，反映节点所处对话情境的传播张力。二者的预测结果在逻辑上可能出现分离——一条评论可以在内容层面具有较高传播价值（高回复吸引力），但处于已趋于收敛的讨论链末端（低路径延续概率）；反之，一条内容普通的评论也可能因所处交互序列的动态特征而具有较高路径延续概率。本模块将两个维度联合呈现，避免单一指标遮蔽传播特征的某一侧面，在"评论质量"与"讨论情境"的交叉视角下更完整地刻画多主体的传播行为差异。

\textbf{（1）分层聚合}

系统按主体类型将评估结果分为Human子集和Agent子集，分别统计两个子集在两个评估维度上的表现。

\textbf{（2）传播价值维度对比}

基于评论传播价值评估模块的输出，该步骤计算Human与Agent评论传播价值评分的分布对比（均值、中位数、分位数），统计不同评分区间内Human和Agent的占比分布，并识别高价值评论（评分大于阈值）中Human和Agent的比例。

\textbf{（3）传播延续维度对比}

基于传播延续预测模块的输出，该步骤计算Human与Agent节点2跳增长概率的分布对比，分析不同主体类型节点引发持续讨论的概率差异，并对比含不同交互类型（H$\to$H、H$\to$A、A$\to$H、A$\to$A）的路径在传播延续上的表现差异。

\textbf{（4）情感倾向维度对比}

该步骤从CommentFeature表读取正面、中性、负面三类情感概率，按主体类型分组，计算Human与Agent评论各情感概率的分布均值对比。从内容认知特征角度呈现两类主体在内容表达偏向（情感倾向）上的系统性差异，与传播价值和传播延续维度形成互补：两类主体的行为差异（是否获得回复、是否引发延续）可与内容特征差异（情感偏向）对照，为解读传播机制差异提供认知层面的参照。

\textbf{（5）综合传播特征报告生成}

本模块将三个维度的分层统计结果（Human/Agent各自的传播价值评分分布参数、2跳增长概率分布参数、情感概率均值、各交互类型传播延续均值）结构化为固定格式的提示词（prompt），调用配置的大语言模型后端（可配置为本地推理或OpenAI-compatible API），生成一段自然语言的多主体传播特征分析报告。报告结构固定为三个段落：（一）两类主体传播价值分布差异及解读；（二）两类主体传播延续能力差异及交互类型影响；（三）综合结论，对当前数据集下的多主体信息认知传播特征作出简要判断。LLM的输入仅为聚合统计量（均值、中位数、分位数等数值），不包含原始评论文本；提示词采用固定模板将统计量填入结构化指令，保证输出格式的可复现性。三个维度的聚合统计数据与LLM生成报告一并写入AnalysisReport表持久化存储，供仪表盘展示和导出使用。

多主体对比分析模块的交互时序如图\ref{fig:analysis_seq}所示，展示了从触发分析到报告生成的完整协作流程。API路由层首先校验导入任务是否存在且评估结果是否已生成，随后将请求委派至analysis\_service模块的run\_analysis()函数；该函数联表查询Comment、CommentFeature与EvaluationResult后，按主体类型将数据分为Human子集和Agent子集，依次调用\_compute\_group\_stats()计算两个子集在传播价值、传播延续和情感倾向三个维度的分布统计量（均值、中位数、四分位数），再调用\_compute\_interaction\_stats()根据评论树父子关系动态计算各交互类型（H$\to$H / H$\to$A / A$\to$H / A$\to$A）下的2跳增长概率均值。最后调用\_generate\_llm\_report()将统计结果填入固定提示词模板并通过OpenAI-compatible API请求LLM后端生成三段式自然语言分析报告，统计数据与报告文本一并写入AnalysisReport表。

\begin{figure}[H]
  \centering
  \includegraphics[width=0.9\textwidth]{G6-4/图6-4 多主体对比分析模块时序图.cropped.pdf}
  \bicaption{多主体对比分析模块时序图}{Sequence Diagram of Multi-agent Comparative Analysis Module}
  \label{fig:analysis_seq}
\end{figure}

\subsection{评估结果展示与管理模块的实现}

该模块对应第三章设计的第六个功能模块，整合评估结果查询、数据集管理与可视化交互功能，提供统一的Web界面。

\textbf{（1）评估结果查询}

该功能支持按主体类型（Human/Agent）、传播价值评分区间、2跳增长概率区间对评估结果进行组合检索，同时支持按话题关键词或地区筛选评论树。查询结果以列表形式呈现，支持按评分或概率排序，可从结果条目直接跳转至评论树可视化视图。

\textbf{（2）数据集管理}

该功能提供所有已导入数据集的元信息查看（文件名、导入时间、各阶段样本量、评估状态），支持删除导入任务及其关联的评论数据、特征数据与评估结果，并支持对已导入数据集重新触发评估推理（如模型权重更新后对原始数据重新评分）。

\textbf{（3）评论树可视化}

基于D3.js\cite{bostock2011d3}实现，以树形结构展示评论树，节点标注主体类型（不同颜色区分Human与Agent），并在节点上叠加传播价值评分（颜色深浅表示评分高低）。用户可点击节点查看评论详情，包括文本内容、评分和预测概率等信息。

\textbf{（4）路径高亮与探索}

该功能高亮显示高2跳增长概率的讨论路径，在路径上标注交互类型（H$\to$H、H$\to$A、A$\to$H、A$\to$A），支持按主体类型和预测概率区间进行路径筛选。

\textbf{（5）多主体对比仪表盘}

基于ECharts实现，包含传播价值评分分布对比图（Human与Agent评分分布曲线或箱型图）、2跳增长概率分布对比图（Human与Agent概率分布曲线）、情感倾向分布对比图（以分组柱状图展示Human与Agent的正、中、负情感概率均值）、分交互类型传播延续对比图（以H$\to$H、H$\to$A、A$\to$H、A$\to$A为分组对比各交互类型下2跳增长概率均值），以及关键指标对比卡片（分别展示Human与Agent的传播价值评分均值、2跳增长概率均值和正面情感概率均值）。

\textbf{（6）数据导出}

该功能支持将多主体对比统计数据和评估结果列表导出为CSV格式。

评估结果展示与管理模块的核心交互时序如图\ref{fig:display_seq}所示，以"评估结果查询后跳转评论树可视化"场景为例展示前后端交互流程。用户在查询面板设置筛选条件（主体类型、评分区间、概率区间等）后，前端向GET /api/comments接口发送请求，QueryRouter通过Comment、EvaluationResult与CommentFeature三表左连接查询并应用筛选、排序与分页条件，返回评论列表；用户点击某条结果后，前端携带根节点ID请求GET /api/tree/\{root\_tweet\_id\}接口，TreeRouter加载该评论树全部节点及对应的评估结果，通过\_build\_tree()递归构建嵌套的树节点结构（每个节点包含评论ID、主体类型、传播价值评分、增长概率及子节点列表）返回前端，前端基于D3.js渲染树形可视化视图，完成从列表到图表的跳转交互。

\begin{figure}[H]
  \centering
  \includegraphics[width=0.9\textwidth]{G6-5/图6-5 评估结果展示与管理模块交互时序图.cropped.pdf}
  \bicaption{评估结果展示与管理模块交互时序图}{Interaction Sequence Diagram of Result Display and Management Module}
  \label{fig:display_seq}
\end{figure}

\section{系统测试}

\subsection{测试环境}

系统测试在表\ref{tab:test_env}所列的软硬件环境下进行。

\begin{table}[htbp]
  \centering
  \bicaption{测试环境配置}{Test Environment Configuration}
  \label{tab:test_env}
  \renewcommand\arraystretch{1.5}
  \begin{tabular}{>{\centering\arraybackslash}m{2cm} >{\centering\arraybackslash}m{3.5cm} >{\centering\arraybackslash}m{7cm}}
    \toprule
    类\quad 别 & 项\quad 目 & 配\quad 置 \\
    \midrule
    硬件 & 处理器 & Apple M4（10核：4性能核 + 6能效核） \\
    硬件 & 内存 & 24 GB 统一内存 \\
    硬件 & 存储 & 512 GB SSD \\
    操作系统 & 系统版本 & macOS 15.3.2（Darwin arm64） \\
    运行时 & Python & 3.13.9 \\
    运行时 & Node.js & 25.1.0 \\
    后端框架 & FastAPI & 0.133.1 \\
    后端框架 & SQLAlchemy & 2.0.47 \\
    深度学习 & PyTorch & 2.9.1（MPS加速） \\
    深度学习 & SentenceTransformers & 5.2.0 \\
    测试框架 & Playwright & 1.58.2（Chromium） \\
    数据库 & SQLite & 3（嵌入式，随Python发行） \\
    \bottomrule
  \end{tabular}
\end{table}

\subsection{功能测试}

针对系统各功能模块，设计功能测试用例并进行验证。

\subsubsection{数据导入与预处理功能测试}

数据导入与存储验证测试用例如表\ref{tab:test_import}所示，数据导入功能效果如图\ref{fig:test_import}所示。

\begin{table}[htbp]
  \centering
  \bicaption{数据导入与存储验证测试用例}{Test Case for Data Import and Storage Verification}
  \label{tab:test_import}
  \renewcommand\arraystretch{1.3}
  \footnotesize
  \begin{tabular}{>{\centering\arraybackslash}m{0.8cm} >{\raggedright\arraybackslash}m{5.7cm} >{\raggedright\arraybackslash}m{6cm}}
    \toprule
    序号 & \centering 输入及操作说明 & \centering\arraybackslash 期望测试结果 \\
    \midrule
    1 & 准备包含评论树原始数据的JSON文件，文件为JSON数组，每条记录包含11个必要字段，数据中包含重复评论ID记录、采集时间缺失记录及随访时间不足12小时的记录 & JSON文件格式合法，包含可触发去重、质量过滤和右删失过滤的多类测试数据 \\
    2 & 进入"数据管理"页面，点击上传区域选择该JSON文件上传 & 页面显示蓝色提示"正在处理，请稍候……" \\
    3 & 等待系统完成预处理流水线（格式校验、去重、质量过滤、右删失过滤、评论树重构、特征提取、数据库写入） & 页面弹出成功提示，显示各阶段计数逐步递减 \\
    4 & 点击"刷新"按钮，在数据集列表中查看新增记录 & 列表新增一行，文件名与上传文件一致，各计数字段与成功提示一致，状态显示绿色"completed" \\
    \midrule
    \multicolumn{3}{l}{测试结果：通过\hspace{2em}异常：无} \\
    \bottomrule
  \end{tabular}
\end{table}

\begin{figure}[H]
  \centering
  \includegraphics[width=0.9\textwidth]{TEST/数据导入与存储验证测试用例.png}
  \bicaption{数据导入功能效果图}{Screenshot of Data Import Function}
  \label{fig:test_import}
\end{figure}

异常数据导入测试用例如表\ref{tab:test_import_err}所示，异常数据导入提示效果如图\ref{fig:test_import_err}所示。

\begin{table}[htbp]
  \centering
  \bicaption{异常数据导入测试用例}{Test Case for Abnormal Data Import}
  \label{tab:test_import_err}
  \renewcommand\arraystretch{1.3}
  \footnotesize
  \begin{tabular}{>{\centering\arraybackslash}m{0.8cm} >{\raggedright\arraybackslash}m{5.7cm} >{\raggedright\arraybackslash}m{6cm}}
    \toprule
    序号 & \centering 输入及操作说明 & \centering\arraybackslash 期望测试结果 \\
    \midrule
    1 & 准备一个缺少必要字段的JSON文件，例如第1条记录中缺少tweet\_id和created\_at字段 & JSON文件已准备，包含字段不完整的记录 \\
    2 & 进入"数据管理"页面，点击上传区域上传该异常JSON文件 & 系统执行格式校验 \\
    3 & 观察系统返回的错误提示 & 页面显示红色错误提示，明确指出缺失字段名称，数据未写入数据库 \\
    \midrule
    \multicolumn{3}{l}{测试结果：通过\hspace{2em}异常：无} \\
    \bottomrule
  \end{tabular}
\end{table}

\begin{figure}[H]
  \centering
  \includegraphics[width=0.9\textwidth]{TEST/异常数据导入测试用例.png}
  \bicaption{异常数据导入提示效果图}{Screenshot of Abnormal Data Import Error Prompt}
  \label{fig:test_import_err}
\end{figure}

\subsubsection{评论传播价值评估功能测试}

评论传播价值评估功能测试用例如表\ref{tab:test_value}所示，评估结果效果如图\ref{fig:test_value}所示。

\begin{table}[htbp]
  \centering
  \bicaption{评论传播价值评估功能测试用例}{Test Case for Comment Propagation Value Assessment}
  \label{tab:test_value}
  \renewcommand\arraystretch{1.3}
  \footnotesize
  \begin{tabular}{>{\centering\arraybackslash}m{0.8cm} >{\raggedright\arraybackslash}m{5.7cm} >{\raggedright\arraybackslash}m{6cm}}
    \toprule
    序号 & \centering 输入及操作说明 & \centering\arraybackslash 期望测试结果 \\
    \midrule
    1 & 在"数据管理"页面选择状态为"completed"的数据集，点击操作列的绿色"评估"按钮 & 系统调用Group-aware Fusion TabTransformer和IHPE模型依次执行批量推理，弹出提示显示评估总数、传播价值评分数和增长预测数 \\
    2 & 切换至"评估结果"页面，选择该数据集，排序选"传播价值$\downarrow$"，点击"查询" & 结果列表按传播价值评分降序排列，每行显示主体类型、评论内容、传播价值评分（带颜色条）、增长概率、情感概率、发布时间 \\
    3 & 观察"传播价值"列的评分值与颜色条 & 所有评分在[0, 1]区间内，颜色条：$\geq$0.7绿色、$\geq$0.4橙色、$<$0.4红色 \\
    4 & 设置"传播价值$\geq$"为0.5，点击"查询" & 返回结果中所有评论的传播价值评分均$\geq$0.5 \\
    \midrule
    \multicolumn{3}{l}{测试结果：通过\hspace{2em}异常：无} \\
    \bottomrule
  \end{tabular}
\end{table}

\begin{figure}[H]
  \centering
  \includegraphics[width=0.9\textwidth]{TEST/评论传播价值评估测试用例.png}
  \bicaption{评论传播价值评估结果效果图}{Screenshot of Comment Propagation Value Assessment Results}
  \label{fig:test_value}
\end{figure}

\subsubsection{传播延续预测功能测试}

传播延续预测功能测试用例如表\ref{tab:test_growth}所示，预测结果效果如图\ref{fig:test_growth}所示。

\begin{table}[htbp]
  \centering
  \bicaption{传播延续预测功能测试用例}{Test Case for Propagation Continuation Prediction}
  \label{tab:test_growth}
  \renewcommand\arraystretch{1.3}
  \footnotesize
  \begin{tabular}{>{\centering\arraybackslash}m{0.8cm} >{\raggedright\arraybackslash}m{5.7cm} >{\raggedright\arraybackslash}m{6cm}}
    \toprule
    序号 & \centering 输入及操作说明 & \centering\arraybackslash 期望测试结果 \\
    \midrule
    1 & 在"评估结果"页面，选择已完成评估的数据集，排序切换为"增长概率$\downarrow$"，点击"查询" & 结果列表按2跳增长概率从高到低排列 \\
    2 & 观察"增长概率"列的概率值与颜色条 & 所有节点的增长概率在[0, 1]区间内，颜色条按阈值正确显示 \\
    3 & 设置"增长概率$\geq$"为0.5，排序切换为"树深度$\downarrow$"，点击"查询" & 返回结果中所有节点增长概率均$\geq$0.5，按树深度降序排列 \\
    4 & 清除筛选条件，分别观察树深度较小（$\leq$2）和较大（$\geq$5）的节点的增长概率 & 短路径和长路径节点均有有效概率值（非空、在[0, 1]区间内），IHPE模型对不同长度路径均正常输出 \\
    \midrule
    \multicolumn{3}{l}{测试结果：通过\hspace{2em}异常：无} \\
    \bottomrule
  \end{tabular}
\end{table}

\begin{figure}[H]
  \centering
  \includegraphics[width=0.9\textwidth]{TEST/传播延续预测测试用例.png}
  \bicaption{传播延续预测结果效果图}{Screenshot of Propagation Continuation Prediction Results}
  \label{fig:test_growth}
\end{figure}

\subsubsection{多主体对比分析功能测试}

多主体对比分析与报告生成测试用例如表\ref{tab:test_analysis}所示，对比分析效果如图\ref{fig:test_analysis}所示。

\begin{table}[htbp]
  \centering
  \bicaption{多主体对比分析与报告生成测试用例}{Test Case for Multi-agent Comparative Analysis and Report Generation}
  \label{tab:test_analysis}
  \renewcommand\arraystretch{1.3}
  \footnotesize
  \begin{tabular}{>{\centering\arraybackslash}m{0.8cm} >{\raggedright\arraybackslash}m{5.7cm} >{\raggedright\arraybackslash}m{6cm}}
    \toprule
    序号 & \centering 输入及操作说明 & \centering\arraybackslash 期望测试结果 \\
    \midrule
    1 & 确认数据集已完成评估，在"数据管理"页面点击该数据集操作列的蓝色"分析"按钮 & 系统按主体类型分层计算三个维度的统计量，调用LLM后端生成分析报告，弹出成功提示显示两类主体的评论数量 \\
    2 & 切换至"对比分析"页面，选择该数据集，点击"加载"按钮 & 页面显示8张关键指标卡片：Human/Agent评论数、双方传播价值评分均值、双方增长概率均值、双方正面情感概率均值 \\
    3 & 查看传播价值评分分布与2跳增长概率分布对比图 & 分组柱状图以10个区间为横轴，蓝色系列Human、橙色系列Agent，各区间频次之和等于对应主体评论数 \\
    4 & 查看情感倾向对比图和交互类型$\times$增长概率均值图 & 情感对比图横轴为正面、中性、负面三类；交互类型图横轴为H$\to$H、H$\to$A、A$\to$H、A$\to$A \\
    5 & 查看"LLM多主体传播行为分析报告"区域，点击"导出文本"按钮 & 报告包含三段内容：传播价值差异、传播延续差异与交互类型影响、综合结论；点击导出后浏览器下载报告文本文件 \\
    \midrule
    \multicolumn{3}{l}{测试结果：通过\hspace{2em}异常：无} \\
    \bottomrule
  \end{tabular}
\end{table}

\begin{figure}[H]
  \centering
  \includegraphics[width=0.9\textwidth]{TEST/多主体对比分析与报告生成测试用例.png}
  \bicaption{多主体对比分析与报告生成效果图}{Screenshot of Multi-agent Comparative Analysis and Report Generation}
  \label{fig:test_analysis}
\end{figure}

\subsubsection{评估结果展示与管理功能测试}

数据集管理功能测试用例如表\ref{tab:test_manage}所示，管理效果如图\ref{fig:test_manage}所示。

\begin{table}[htbp]
  \centering
  \bicaption{数据集管理功能测试用例}{Test Case for Dataset Management}
  \label{tab:test_manage}
  \renewcommand\arraystretch{1.3}
  \footnotesize
  \begin{tabular}{>{\centering\arraybackslash}m{0.8cm} >{\raggedright\arraybackslash}m{5.7cm} >{\raggedright\arraybackslash}m{6cm}}
    \toprule
    序号 & \centering 输入及操作说明 & \centering\arraybackslash 期望测试结果 \\
    \midrule
    1 & 进入"数据管理"页面，查看数据集列表 & 列表显示各数据集元信息：ID、文件名、导入时间、各阶段样本量、评论树数、状态 \\
    2 & 点击某"completed"数据集操作列的红色"删除"按钮 & 弹出确认对话框，提示此操作将级联删除所有相关数据 \\
    3 & 在确认对话框中点击"确定" & 数据集从列表移除，关联的评论、特征、评估结果和分析报告均被级联删除 \\
    4 & 点击"刷新"刷新列表 & 被删除数据集不再出现在列表中 \\
    \midrule
    \multicolumn{3}{l}{测试结果：通过\hspace{2em}异常：无} \\
    \bottomrule
  \end{tabular}
\end{table}

\begin{figure}[H]
  \centering
  \includegraphics[width=0.9\textwidth]{TEST/数据集管理功能测试用例.png}
  \bicaption{数据集管理效果图}{Screenshot of Dataset Management}
  \label{fig:test_manage}
\end{figure}

评论树节点传播价值与边增长概率展示测试用例如表\ref{tab:test_tree}所示，展示效果如图\ref{fig:test_tree}所示。

\begin{table}[htbp]
  \centering
  \bicaption{评论树节点传播价值与边增长概率展示测试用例}{Test Case for Comment Tree Node Propagation Value and Edge Growth Probability Display}
  \label{tab:test_tree}
  \renewcommand\arraystretch{1.3}
  \footnotesize
  \begin{tabular}{>{\centering\arraybackslash}m{0.8cm} >{\raggedright\arraybackslash}m{5.7cm} >{\raggedright\arraybackslash}m{6cm}}
    \toprule
    序号 & \centering 输入及操作说明 & \centering\arraybackslash 期望测试结果 \\
    \midrule
    1 & 进入"评论树"页面，选择已评估数据集，在下拉框中选择一棵节点较多的评论树，点击"重绘" & 页面渲染树形结构，Human节点蓝色、Agent节点橙色，节点半径与传播价值评分成正比 \\
    2 & 观察树中各节点的大小差异 & 高传播价值节点圆形明显更大、颜色更饱满，低传播价值节点圆形较小、颜色较淡 \\
    3 & 观察树中连边的颜色与粗细差异 & 增长概率$\geq$阈值的边为红色粗线，$\geq$阈值$\times$0.6的边为橙色中等线，其余为浅灰细线 \\
    4 & 调整"高亮阈值（增长概率）"数值后点击"重绘" & 阈值提高后红色高亮边减少，阈值降低后红色高亮边增多 \\
    5 & 将鼠标悬停在某节点上，查看浮动提示信息 & 显示用户名、主体类型、传播价值评分、2跳增长概率、评论文本前120字符 \\
    \midrule
    \multicolumn{3}{l}{测试结果：通过\hspace{2em}异常：无} \\
    \bottomrule
  \end{tabular}
\end{table}

\begin{figure}[H]
  \centering
  \includegraphics[width=0.9\textwidth]{TEST/评论树节点传播价值与边增长概率展示测试用例.png}
  \bicaption{评论树节点传播价值与边增长概率展示效果图}{Screenshot of Comment Tree Visualization}
  \label{fig:test_tree}
\end{figure}

\subsection{性能测试}

使用与功能测试相同的数据集（356条评论，1棵评论树）开展性能测试，通过Playwright记录各操作的端到端耗时。数据导入与预处理流水线整体耗时约5.4秒，主要开销来自BGE文本编码与情感分析两个深度学习特征提取步骤；双模型推理（Group-aware Fusion TabTransformer与IHPE）耗时约1.3秒，得益于模型在应用启动时已预加载且参数规模较小；多主体对比分析耗时约56.1秒，瓶颈在于调用本地LLM后端生成分析报告，统计聚合本身不足1秒。前端各交互操作（评估结果查询、评论树可视化渲染、对比仪表盘加载）均在1.2秒以内完成。综上，系统在中小规模数据集下各环节响应时间均处于可接受范围，满足交互式评估分析的使用需求。

\section{本章小结}

本章在第三章系统设计的基础上，阐述了多主体信息认知传播评估系统六个功能模块从数据预处理、存储、模型推理到对比分析与可视化展示的具体实现方案。功能测试与性能测试结果表明，系统端到端流程正确，各环节响应效率满足交互式分析的使用需求。
